% !TeX root = ../main.tex

\chapter{中國文學}

\section{出師表}

\begingroup
\centering

\begin{tabular}{ll}
    What & is     \\
    this & doing? \\
\end{tabular}
\captionsetup{type=table}
\captionof{table}{This is an example table.}\label{tbl:nicetablelesstable}
\endgroup


\begin{figure}[htbp]
    \centerline{\includegraphics[width=0.5\columnwidth]{seal.png}}
    \caption{Example of a figure caption.}
    \label{fig}
\end{figure}

Please use `\textbackslash eqref' to refer an equation. Use `\textbackslash ref' to refer sections, figures or tables. For example, this is Section \ref{sec:1}.

\ncsred{臣亮言}：先帝創業未半，而中道崩殂。今天下三分，益州疲弊，此誠危急存亡之秋也。然侍衛之臣，不懈於內；忠志之士，忘身於外者，蓋追先帝之殊遇，欲報之於陛下也。誠宜開張聖聽，以光先帝遺德，恢弘志士之氣；不宜妄自菲薄，引喻失義，以塞忠諫之路也。宮中府中，俱為一體，陟罰臧否，不宜異同。若有作姦犯科，及為忠善者，宜付有司，論其刑賞，以昭陛下平明之治，不宜篇私，使內外異法也。\par

侍中、侍郎郭攸之、費褘、董允等，此皆良實，志慮忠純，是以先帝簡拔以遺陛下。愚以為宮中之事，事無大小，悉以咨之，然後施行，必能裨補闕漏，有所廣益。將軍向寵，性行淑均，曉暢軍事，試用於昔日，先帝稱之曰「能」，是以眾議舉寵為督。愚以為營中之事，悉以咨之，必能使行陣和睦，優劣得所。親賢臣，遠小人，此先漢所以興隆也；親小人，遠賢臣，此後漢所以傾頹也。先帝在時，每與臣論此事，未嘗不歎息痛恨於桓、靈也。侍中、尚書、長史；參軍，此悉貞良死節之臣也，願陛下親之信之，則漢室之隆，可計日而待也。

臣本布衣，躬耕於南陽，苟全性命於亂世，不求聞達於諸侯。先帝不以臣卑鄙，猥自枉屈，三顧臣於草廬之中，諮臣以當世之事，由是感激，遂許先帝以驅馳。後值傾覆，受任於敗軍之際，奉命於危難之間，爾來二十有一年矣！先帝知臣謹慎，故臨崩寄臣以大事也。受命以來，夙夜憂勤，恐託付不效，以傷先帝之明。故五月渡瀘，深入不毛。今南方已定，兵甲已足，當獎率三軍，北定中原，庶竭駑鈍，攘除奸凶，興復漢室，還於舊都；此臣所以報先帝而忠陛下之職分也。至於斟酌損益，進盡忠言，則攸之、褘、允之任也。

願陛下託臣以討賊興復之效；不效，則治臣之罪，以告先帝之靈。若無興德之言，則戮允等，以彰其慢。陛下亦宜自課，以諮諏善道，察納雅言，深追先帝遺詔，臣不勝受恩感激。

今當遠離，臨表涕泣，不知所云。

\section{短歌行\label{sec:1}}

對酒當歌，人生幾何！譬如朝露，去日苦多。
慨當以慷，憂思難忘。何以解憂？唯有杜康。
青青子衿，悠悠我心。但為君故，沉吟至今。
呦呦鹿鳴，食野之苹。我有嘉賓，鼓瑟吹笙。
明明如月，何時可掇？憂從中來，不可斷絕。
越陌度阡，枉用相存。契闊談宴，心念舊恩。
月明星稀，烏鵲南飛。繞樹三匝，何枝可依？
山不厭高，海不厭深。周公吐哺，天下歸心。

\section{始得西山宴遊記}

自余爲僇人，居是州，恆惴慄。其隙也，則施施而行，漫漫而遊。日與其徒上高山，入深林，窮回溪；幽泉怪石，無遠不到。到則披草而坐，傾壺而醉，醉則更相枕以臥，臥而夢。意有所極，夢亦同趣。覺而起，起而歸。以爲凡是州之山有異態者，皆我有也，而未始知西山之怪特。

今年九月二十八日，因坐法華西亭，望西山，始指異之。遂命僕人過湘江，緣染溪，斫榛莽，焚茅茷，窮山之高而止。攀援而登，箕踞而遨，則凡數州之土壤，皆在衽席之下

其高下之勢，岈然窪然，若垤若穴，尺寸千里，攢蹙累積，莫得遁隱；縈青繚白，外與天際，四望如一。然後知是山之特出，不與培塿爲類。悠悠乎與灝氣俱，而莫得其涯；洋洋乎與造物者遊，而不知其所窮。

引觴滿酌，頹然就醉，不知日之入，蒼然暮色，自遠而至，至無所見，而猶不欲歸。心凝形釋，與萬化冥合。然後知吾向之未始遊，遊於是乎始，故爲之文以志。是歲，元和四年也。